\documentclass[11pt]{article}
\usepackage[margin=1in]{geometry}
\usepackage[T1]{fontenc}
\usepackage{lmodern,amsmath,booktabs,graphicx}
\usepackage[hidelinks]{hyperref}
\usepackage{xurl}
\setlength{\parindent}{0pt}\setlength{\parskip}{6pt}
\setlength{\emergencystretch}{3em}
\title{Causal NRMM Body-Heading Feedback for LiDAR Rectangle Localization}
\author{Pan Dynamics Collision Avoidance Research}
\date{October 3, 2026}
\begin{document}\maketitle
\section{Implemented behavior and measured outcome}
The perception pipeline can now receive a predicted target body heading from
an estimator, fit a rectangle with that orientation held fixed, and return the
new target position to the estimator. The current MATLAB NRMM runtime already
reconstructs this heading. Its state equations, configuration and gains were
not modified. Both initial and inlier-refinement rectangle fits preserve the
supplied orientation, modulo the rectangle's 180-degree symmetry.

Actual YOLO26x inference and native LiDAR association were rerun on two
previously recorded parked-free CARLA datasets. The following comparison keeps
symmetric face placement, declared 5.0-by-1.9 m dimensions, confidence selection,
filters and native ego preprocessing unchanged. The previous method estimates
a motion direction from the position history; the new method uses the causal
NRMM body-heading prediction. Each recording contains 300 frames at 0.02 s.

\begin{center}\scriptsize
\begin{tabular}{@{}llrrrrr@{}}\toprule
 & & \multicolumn{3}{c}{After 2 seconds} & \multicolumn{2}{c}{All frames} \\
Data & Heading source & Raw pos. & Est. pos. & Est. vel. & Est. pos. & Est. vel. \\
 & & m & m & m/s & m & m/s \\\midrule
15 m & Previous motion heading & 0.0674 & 0.0946 & 0.3049 & 0.1314 & 0.8195 \\
15 m & NRMM body-heading feedback & 0.0810 & 0.1274 & 0.4471 & 0.1557 & 0.8663 \\
20 m & Previous motion heading & 0.0661 & 0.1430 & 0.4997 & 0.6948 & 2.6381 \\
20 m & NRMM body-heading feedback & 0.0719 & 0.1018 & 0.3782 & 0.1476 & 0.8506 \\
\bottomrule\end{tabular}\end{center}

The 20 m recording's cold-start behavior improves substantially, but the
15 m recording worsens relative to the preceding symmetric-placement method.
This is not a uniform accuracy improvement, and the optional feedback path
has not replaced the default. No new scene or independent capture was used.
All-frame and post-acquisition metrics are both retained; estimator RMSE
includes dropout intervals while raw position RMSE uses available detections.

\section{Which heading is used}
The existing six-state target chain contains relative position $\hat\rho$,
target absolute velocity $\hat q$ expressed in ego body axes, and target
acceleration $\hat s$ in those axes. The reconstructed body orientation is
\[
 \hat\psi_{C/E}=\operatorname{atan2}(\hat q_y,\hat q_x)-\hat\beta_C.
\]
The implementation exposes it as \path{targetRelativeHeading}; it also
publishes \path{targetHeadingInertial}. The fitted rectangle needs body
heading, including the estimator's sideslip correction, not just velocity
course. Since the predicted NRMM state is already in the current ego body
axes, the native CARLA adapter negates \path{targetRelativeHeading} to match
CARLA's opposite planar angular convention. No world-yaw estimate is needed.
Converting through estimated inertial heading and subtracting the native ego
yaw would introduce an extra ego-heading estimation error. Earlier attempts
with that adapter are preserved separately as diagnostic artifacts and are
excluded from the final tables.

The generic Python contract is \path{HeadingPrior}, with fields
\path{heading_ego_rad}, \path{state_time}, \path{last_measurement_time}, and
\path{source}. The caller supplies it in the
same current ego axes as the point cloud. Zero-speed direction is not
observable from the velocity state; an adapter can return no prior, in which
case the existing geometric search remains available.

\section{Causal ordering and initialization}
The implemented ordering at sample $k$ is:
\begin{enumerate}
\item The previous \texttt{sample} call has integrated measurements through
$k-1$ to a state predicted at $t_k$.
\item Query the current predicted body heading without advancing or modifying
that state. The prior's state time must equal $t_k$, and its latest measurement
time must be strictly earlier than $t_k$.
\item Fit the current detected LiDAR cluster once at the supplied orientation,
including the inlier refinement. No same-frame target-position update has yet
been applied, and no motion-heading iteration revises this result.
\item Supply the fitted position, or an explicit missing measurement, to
\path{onlineNrmmTrackingRuntime} with the \texttt{sample} action. This advances the observer
for the next frame. Previously published positions are never revisited.
\end{enumerate}

The first frame necessarily has no converged target-heading estimate.
This replay uses the existing co-moving initialization: native GNSS ego
velocity initializes target velocity, target acceleration is zero, and the
first fitted position initializes relative position. The first heading is
explicitly labeled \texttt{initialization}; later ones are \texttt{prediction}.
That prior is suitable for these following-vehicle scenes but is an assumption,
not a heading measurement or a solution to arbitrary target acquisition.

The API is invoked through \path{run_pipeline(args,heading_feedback=adapter)}.
An adapter implements \path{predict(frame,time)}
and \path{update(frame,time,position)}. This keeps estimator transport outside
the geometry implementation. The local research adapter uses the actual
installed MATLAB Engine and production runtime; no Python approximation of
NRMM or actor-truth input is substituted. Recorded detections can be cached
because their association is independent of heading; the final runs rerun
actual detection. This remains an offline replay, not a 50 Hz implementation.

\section{Ablations and failure cases}
Cached actual detections and native clouds were used for the diagnostic
ablations. With original edge anchoring, feedback is much worse than with
symmetric placement: post-2-second position/velocity RMSE is
0.5927 m / 2.0621 m/s at 15 m and 0.3377 m / 1.1868 m/s at 20 m.
This shows that supplying heading does not by itself eliminate placement bias.

For the straight 20 m recording, holding the initial world heading constant
instead of updating it gives 0.0256 m position / 0.1101 m/s velocity RMSE
after 2 s. This is a revealing control, not a recommendation for turning
vehicles: a good initial straight-line assumption already supplies much of
the useful information, while noisy feedback can degrade it.

When the initial target velocity/body-heading prior is rotated by 20 degrees,
the 20 m feedback case has 0.6761 m position / 2.5280 m/s velocity RMSE after
2 s. At 90 degrees it has 0.7598 m / 2.6344 m/s. The corresponding full-record
position RMSE values are 0.9618 and 1.3008 m. All cases remain finite, but
several accuracy gates fail. Hard orientation locking can preserve a wrong
initial interpretation rather than reliably recover from it.

The measurement now depends on the observer's prior state: orientation error
changes the fitted center, which subsequently changes velocity and heading
estimates. Deterministic ISS analysis does not require statistical independence,
but the assumed measurement-error and model-domain bounds must still hold for
this feedback path. These tests do not establish those bounds or a new
coupled-system stability certificate. A bounded
heading search, uncertainty-aware acceptance or multi-hypothesis acquisition
would require separate implementation and validation; none is claimed here.

\section{Validation and reproducibility}
All 32 Python geometry, association, causality and feedback behavior tests
pass. They check fixed orientation in both fitting stages, current body-frame
semantics, strict prior timestamps, predict-before-update order, missing
measurements, duplicate-time rejection, nonmutation and prefix invariance.
Four focused MATLAB tests pass: current-publication/next-state equivalence,
read-only output queries, analytic maneuver convergence including target
heading, and pure prediction during radar dropout.

The actual-feedback normal/blackout pair in each recording differs only in
RGB paths at frames 150--164. Audit results record fixed-angle residuals,
strictly past prior timestamps, finite output counts, matching prefixes,
reacquisition and unchanged ego states. Original raw files are preserved.
Scoring hashes are written before the separate scorer reads actor truth.
The initial GPU attempt failed from insufficient shared memory. Bounded-core
CPU attempts were interrupted for the frame correction and, later, when GPU
memory became available. Final inference uses GPU 0 with a 20 percent PyTorch
allocation cap and lower scheduling priority, without touching other processes
or any CARLA server. Complete commands, runtime versions, source
hashes, diagnostics and actual timings accompany the report.

Configured target speeds start at 10 m/s, while these recordings include
slower targets. Some raw errors exceed the configured 0.12 m measurement
allowance. Availability of an observer-bound flag does not validate those
premises. Native compass heading also retains the earlier ideal-heading
limitation. The generic callback integration and behavior tests are tracked;
raw recordings, local MATLAB replay transport, caches, model weights and
figures remain outside Git, with identified external archive exports.

\IfFileExists{feedback_trace.pdf}{\begin{figure}[p]\centering
\includegraphics[width=\textwidth]{feedback_trace.pdf}
\caption{Actual-feedback errors compared with the previous symmetric geometry.
The vertical line marks the fixed 2 s acquisition exclusion. The shaded band
applies only to the blackout feedback case.}\end{figure}}{}
\end{document}
